% !TeX root = ../main.tex
\section{Conservation of charge}
\label{sec:conservation_of_charge}

Electric charge is carried by the same phase/component material measures as
mass.  Weighting each component balance by the charge carried per unit mass and
summing gives the charge balance.  The field equation for the electric
potential is introduced with the electrical variation; the electrochemical
component potentials are identified later by the Coleman--Noll procedure.

Let \(z_\xi^\alpha\) denote the charge carried per unit mass of component
\(\alpha\) in phase \(\xi\).  The charge density per unit current mixture volume
associated with that component is
%
\begin{equation}
	\varrho_\xi^\alpha
	=
	z_\xi^\alpha\rho_\xi^\alpha
	=
	z_\xi^\alpha\phi_\xi\bar\rho_\xi\eta_\xi^\alpha .
	\label{eq:component_charge_density_definition}
\end{equation}
%
The coefficient \(z_\xi^\alpha\) may vanish for neutral components, may be a
constant valence-to-mass ratio for a mobile ion, or may be an effective charge
per mass for a solid component representing fixed surface or exchange-site
charge.  Fixed surface or exchange-site charge and mobile ionic charge are
represented by the same phase/component variables
\cite{leroy2004triple,newman2004electrochemical}.  A charged component may
belong to any phase; components treated as electrically neutral have
\(z_\xi^\alpha=0\).  The total free charge density is
%
\begin{equation}
	\varrho
	=
	\sum_\xi\sum_\alpha
	z_\xi^\alpha\phi_\xi\bar\rho_\xi\eta_\xi^\alpha.
	\label{eq:mixture_charge_density_definition}
\end{equation}
%
Thus every free-charge contribution is carried by a component mass measure.
The charge density \(\varrho\) is the field-source coefficient in the
electrical virtual work and in Gauss' law, and it multiplies the rate of the
electric potential in the bulk field-source power.  In particular, representing
fixed charge as a solid component ties charge storage, deformation, reaction,
and ion exchange to the same material measures used for mass.

The material charge measure follows by multiplying the augmented component
material measure \eqref{eq:augmented_component_material_measure} by
\(z_\xi^\alpha\),
%
\begin{equation}
	z_\xi^\alpha\rho_{\xi0}^\alpha
	+
	z_\xi^\alpha\mc C_\xi^\alpha
	=
	J_\xi\varrho_\xi^\alpha .
	\label{eq:augmented_component_material_charge_measure}
\end{equation}
%
For constant \(z_\xi^\alpha\), its material rate gives
%
\begin{align}
	\frac{\partial\varrho_\xi^\alpha}{\partial t}
	+
	\nabla_{\mbf x}\cdot
	\left(
	\varrho_\xi^\alpha\mbf v_\xi
	\right)
	&=
	\frac{
	z_\xi^\alpha\dot{\mc C}_\xi^\alpha
	}{J_\xi}.
	\label{eq:component_charge_material_rate}
\end{align}
%
Here and below, \(z_\xi^\alpha\) is constant.  A state-dependent effective
charge requires additional material-rate and transport terms.

Combining \eqref{eq:component_charge_material_rate} with
\eqref{eq:reference_conversion_accumulation_rate} gives the charge balance for
each charged component,
%
\begin{align}
	\frac{\partial\varrho_\xi^\alpha}{\partial t}
	+
	\nabla_{\mbf x}\cdot
	\left[
	z_\xi^\alpha
	\left(
	\rho_\xi^\alpha\mbf v_\xi
	+
	\mbf j_{\xi,{\rm disp}}^\alpha
	+
	\mbf j_{\xi,{\rm diff}}^\alpha
	\right)
	\right]
	&=
	z_\xi^\alpha
	\sum_m\nu_{\xi (m)}^\alpha r_{(m)} .
	\label{eq:component_charge_balance}
\end{align}
%
The current density carried by phase motion and relative component transport is
%
\begin{equation}
	\mbf i
	=
	\sum_\xi\sum_\alpha
	z_\xi^\alpha
	\left(
	\rho_\xi^\alpha\mbf v_\xi
	+
	\mbf j_{\xi,{\rm disp}}^\alpha
	+
	\mbf j_{\xi,{\rm diff}}^\alpha
	\right).
	\label{eq:charge_current_density_definition}
\end{equation}
%
Summing \eqref{eq:component_charge_balance} over phases and components gives
%
\begin{align}
	\frac{\partial}{\partial t}
	\left(
	\sum_\xi\sum_\alpha\varrho_\xi^\alpha
	\right)
	+
	\nabla_{\mbf x}\cdot\mbf i
	=
	\sum_m
	\left(
	\sum_\xi\sum_\alpha
	z_\xi^\alpha\nu_{\xi (m)}^\alpha
	\right)
	r_{(m)} .
	\label{eq:mixture_charge_balance_with_sources}
\end{align}
%
Local internal mechanisms must conserve electric charge,
%
\begin{equation}
	\sum_\xi\sum_\alpha
	z_\xi^\alpha\nu_{\xi (m)}^\alpha
	=0,
	\qquad
	\text{for each mechanism }(m),
	\label{eq:mechanism_charge_conservation}
\end{equation}
%
so the total charge balance is
%
\begin{equation}
	\frac{\partial\varrho}{\partial t}
	+
	\nabla_{\mbf x}\cdot\mbf i
	= 0.
	\label{eq:mixture_charge_balance}
\end{equation}
%
Equation \eqref{eq:mechanism_charge_conservation} is the charge analogue of the
mass-conservation restriction
\(\sum_\xi\sum_\alpha\nu_{\xi (m)}^\alpha=0\).  As one example, an ion-exchange
mechanism between an aqueous phase and charged solid sites may remove an aqueous
component and a solid-site component while creating the exchanged aqueous and
solid-site components.  This mechanism changes fluid composition and solid-site
occupancy while satisfying both mass conservation and
\eqref{eq:mechanism_charge_conservation}.  More generally, the phase/component
formulation admits exchange among any represented phases and components whose
stoichiometry satisfies these two conservation restrictions.  Thus charge
attached to deforming solid sites is transported by the component current
\(\mbf i\).

Ion-exchange and surface-complexation models can be represented by retaining
charged site components in the surface-bearing phase, imposing
\eqref{eq:mechanism_charge_conservation} on the exchange stoichiometry, and
choosing admissible kinetic or equilibrium closures for the corresponding
mechanism rates.  The electrostatic field equation is Gauss' law, derived in
\sref{virtual_power_derivation}.
